%% IntentDiff-RAG — Elsevier CAS double-column
\documentclass[a4paper,fleqn]{cas-dc}
\usepackage[authoryear,longnamesfirst]{natbib}

\begin{document}
\let\WriteBookmarks\relax
\def\floatpagepagefraction{1}
\def\textpagefraction{.001}

\shorttitle{}
\shortauthors{}

\title[mode=title]{IntentDiff-RAG: Lightweight Detection of Malicious npm Updates Using Small Open-Weight LLMs}

\author[1]{}
\cormark[1]
\ead{}
\credit{}

\affiliation[1]{organization={},
            addressline={},
            city={},
            postcode={},
            country={}}

\cortext[1]{Corresponding author}

\begin{abstract}
hallo
\end{abstract}

\begin{highlights}
\item
\item
\item
\end{highlights}

\begin{keywords}
\sep
\end{keywords}

\maketitle


\section{Introduction}\label{sec:intro}

% Every factual claim in this section has been matched to its original source
% (see docs/references.md). Do not add any figure without a verified source.

Modern software development relies on third-party packages distributed through centralized repositories. In the JavaScript ecosystem, the npm registry receives roughly 15 package updates every minute, and installing a single package may pull in dozens of transitive dependencies \citep{wyss2025evaluating}. Every new version of those packages reaches downstream projects through the same mechanism: a version is published, and is then installed by the projects that depend on it. As a result, trusting a package in practice also means trusting every version that will later be published under that package's name.

It is precisely this trust that attackers exploit. Instead of publishing a new package that victims must first discover, attackers inject malicious code into an update of a package that is already widely used, an attack that \citet{sofaer2024rogueone} call a \emph{rogue update}. A series of recent incidents illustrates the scale of this risk. In September 2025, attackers obtained a maintainer's credentials through social engineering and then updated 18 popular packages, including \texttt{chalk} and \texttt{debug}. The injected payload silently manipulated cryptocurrency transactions, while those packages were collectively downloaded roughly two billion times per week \citep{aws2025chalk}. A week later, the \emph{Shai-Hulud} worm began stealing npm tokens and other credentials from developer environments. Using the stolen tokens, it published infected versions of other packages belonging to the same maintainers, and more than 180 packages were compromised within the first 24 hours \citep{aws2025chalk,unit42shaihulud}. In March 2026, the account of one of the \texttt{axios} maintainers was hijacked and used to publish versions 1.14.1 and 0.30.4. Both releases added the dependency \texttt{plain-crypto-js}, which at install time downloaded and executed a remote access trojan \citep{datadog2026axios}. In all three cases, the malicious payload arrived through the legitimate distribution channel of an already trusted package.

The security mechanisms developers commonly rely on are not designed for attacks of this kind. The \texttt{npm audit} command, for instance, sends a description of a project's dependencies to the registry and requests a report of known vulnerabilities \citep{npmaudit}, yet a newly published malicious update has no such record. Detection must therefore be performed on the content of the update itself, before any public report exists. The difficulty is that malicious and benign code use the same APIs. Based on a large-scale evaluation of npm malicious-package detectors, \citet{guo2026effective} conclude that every detector must resolve the ambiguity between the \emph{capability} of code and the \emph{intent} behind it. For example, legitimate build tools such as \texttt{husky}, \texttt{patch-package}, and \texttt{prisma} declare the same install hooks that malware uses. This ambiguity is operationally expensive. Because legitimate updates vastly outnumber malicious ones, even a low false positive rate can produce more false alarms than correct detections
\citep{wyss2025evaluating}.

Prior work has followed two directions. The first is program analysis. RogueOne compares data flows between the old and the new version of a package and flags new data flows that cross a trust domain. On a dataset of single updates, this approach detects 75.6\% of rogue updates with 4.7\% false positives \citep{sofaer2024rogueone}. However, this abstract-interpretation-based analysis requires an average of 226.7 seconds per update, not counting cases that exceed a one-hour timeout. A more sensitive configuration does raise detection to 92.7\%, but at the cost of 40\% false positives \citep{sofaer2024rogueone}. The second direction leverages large language models (LLMs). SocketAI uses GPT-3 and GPT-4 to review npm package code; with GPT-4, the system reaches 99\% precision and 97\% F1 \citep{zahan2025leveraging}. \citet{wyss2025evaluating} focus LLM analysis on package updates. With a smart-diff that reduces input size by 95.4\%, GPT-4o detects all 46 historical malicious updates and 163 malicious clones, with 8 false positives out of 2,000 typical updates. IntelGuard adds retrieval-augmented generation (RAG) with a knowledge base built from 8,024 pages of threat intelligence reports, and then assesses whether the semantics of the code match the package's intended function. On npm packages, the system achieves 98.72\% accuracy with 0.20\% false positives \citep{guo2026bridging}.

Promising as they are, these results leave three open questions. \emph{First}, the success of LLM-based detection has so far been demonstrated with closed models: GPT-4o, GPT-4-Turbo, and Gemini 2.0 Flash in \citet{wyss2025evaluating}, and GPT-4.1-mini in IntelGuard \citep{guo2026bridging}. Closed models require code to be sent to the model provider, and extensive LLM use incurs financial and energy costs that are far from negligible \citep{wyss2025evaluating}. Evaluations of open-weight models that can be run locally do exist, but they were conducted at the level of whole packages and with a model of 6.7 billion parameters \citep{nguyen2025taint}. The ability of models in the 3--4 billion parameter range to assess update smart-diffs has not been measured.

\emph{Second}, retrieval in IntelGuard retrieves only malicious examples \citep{guo2026bridging}. Yet the in-context learning literature shows that the selection and ordering of examples affect a model's predictions. Language models tend to predict labels that appear frequently or that appear near the end of the prompt \citep{zhao2021calibrate}, and examples selected by their similarity to the input yield better performance than random examples \citep{liu2022makes}. A context containing only malicious examples may therefore push the model toward a malicious verdict, including on legitimate updates whose behavior is similar. Positive and negative examples have been proposed as a way to clarify which output characteristics are desired and which the model should avoid \citep{gao2024customizing}. To the best of our knowledge, however, this approach has not been evaluated for malicious package detection. Moreover, the finding that randomly replacing example labels barely degrades performance \citep{min2022rethinking} suggests that contrastive examples must carry explicit reasoning rather than a label alone.

\emph{Third}, the idea of comparing program behavior against the description its developers state is not new. CHABADA clusters Android applications by their descriptions and then flags API usage that deviates from the cluster. Without requiring any known malware pattern, this approach flags 56\% of new malware \citep{gorla2014checking}. IntelGuard applies similar reasoning implicitly \citep{guo2026bridging}. What remains unknown is how much an explicit intent profile, namely the description, README, and scripts from the previous version, can reduce false positives on legitimate updates that add sensitive behavior, when the effect is measured through a controlled ablation.

This study addresses these three questions through IntentDiff-RAG, a static analysis pipeline for assessing npm package updates with small open-weight LLMs. For each version pair, IntentDiff-RAG performs four steps: (i)~constructing a smart-diff following the specification of \citet{wyss2025evaluating}; (ii)~assembling an intent profile from the previous version, since the content of the new version is under the attacker's control \citep{wyss2025evaluating}; (iii)~retrieving one malicious and one legitimate example that are similar in terms of the capabilities the update adds; and (iv)~asking the model to produce a verdict in JSON format together with the lines of code that serve as evidence. The pipeline is evaluated on two models, Qwen3-4B-Instruct-2507 and Qwen2.5-Coder-3B-Instruct. The dataset is split into three subsets: malicious updates, typical updates, and \emph{hard negatives}, that is, legitimate updates that add sensitive capabilities for a documented purpose. Each confirmatory comparison differs in exactly one factor, so that the effect of contrastive examples and the effect of the intent profile can be separated from the effect of the number of examples. This study is guided by the following four research questions.

The contributions of this study are as follows.

\begin{itemize}
\item An evaluation of open-weight LLMs of 3--4 billion parameters for detecting malicious npm updates at the smart-diff level, with metrics reported separately for typical updates and hard negatives.
\item An ablation design that separates the effect of contrastive legitimate examples from the effect of the number of examples, and the effect of the intent profile from the effect of retrieval. The decision rules and the multiple-testing correction were fixed before the experiments were run.
\item A replicable data protocol comprising a doubly labeled hard negative subset, a knowledge card schema, deterministic sample selection, and campaign-level grouping to prevent leakage between the knowledge base and the test data.
\item A measurement of local inference cost, namely tokens, median and 95th
percentile latency, and GPU memory requirements, for each pipeline component.
\end{itemize}

% TODO(results): add one paragraph summarizing the main findings after the
% experiments are complete. Do not fill this in before the result figures are
% available.

The remainder of the article is organized as follows.
Section~\ref{sec:related} discusses related work. Section~\ref{sec:method}
describes IntentDiff-RAG. Section~\ref{sec:setup} presents the data,
configuration, and evaluation procedure. Section~\ref{sec:results} reports the
results, Section~\ref{sec:discussion} discusses the implications and threats to
validity, and Section~\ref{sec:conclusion} concludes the article.

\section{Related Work}\label{sec:related}

% All figures in this section are matched to their original sources (docs/references.md).

This section discusses six lines of research that underpin IntentDiff-RAG:
supply chain attacks on package ecosystems, malicious package detection without
LLMs, malicious update detection, LLM-based detection, the use of retrieval and
examples in prompts for security analysis, and checking the consistency between a
program's description and its behavior. The section closes with the position of
this study relative to these six lines.

\subsection{Supply chain attacks on package ecosystems}

Supply chain risk in npm is rooted in the structure of its ecosystem.
\citet{zimmermann2019small} show that an npm package implicitly trusts an average
of 79 third-party packages and 39 maintainers. They also find that a handful of
maintainer accounts suffices to inject malicious code into the majority of
packages. \citet{ohm2020backstabber} collected and analyzed 174 malicious
packages used in real-world attacks on npm, PyPI, and RubyGems between November
2015 and November 2019. \citet{ladisa2023sok} subsequently developed a taxonomy
of open-source software supply chain attacks in the form of an attack tree. The
taxonomy contains 107 attack vectors linked to 94 real-world incidents and 33
mitigation measures.

This study targets only one class of attack, namely malicious updates to
previously legitimate packages. RogueOne adopts the same restriction, excluding
newly published malicious packages such as typosquats from its scope
\citep{sofaer2024rogueone}. Because many attacks operate at install time, some
research chooses prevention over detection. Latch, for instance, builds a
permission manifest summarizing the behavior of each package at install time and
checks it against a user-defined policy \citep{wyss2022wolf}.

\subsection{Malicious package detection without LLMs}

Early detectors combined several types of analysis. MalOSS combines metadata,
static, and dynamic analysis. Of 339 newly discovered malicious packages reported
to registry maintainers, 278 were confirmed \citep{duan2021towards}. Amalfi
combines a machine learning classifier, a package reproducibility checker, and a
clone detector. Its feature extraction and classification take only a few seconds
per package \citep{sejfia2022practical}.

Later work models behavior as sequences of API calls. DONAPI combines static
analysis with API call sequences from dynamic execution and matches them against
a tiered behavioral knowledge base \citep{huang2024donapi}. Cerebro uses a single
malicious behavior sequence model for both npm and PyPI, and discovered 196 new
malicious packages in npm and 306 in PyPI \citep{zhang2025killing}. OSCAR runs
packages inside a sandbox, fuzzes their exported functions and classes, and
monitors their behavior through API hooks. On npm, OSCAR achieves an F1 of 0.95
and reduces the false positive rate from 34.63\% to 2.57\%
\citep{zheng2024towards}.

The unified evaluation by \citet{guo2026effective} compares 11 tools (16
variants) on 6,420 malicious and 7,288 legitimate packages. From that evaluation,
they conclude that the precision and recall of each tool are determined by how it
resolves the ambiguity between the capability and the intent of the code.

This line of work has two limitations with respect to the problem we study.
First, most of the detectors above assess whole packages rather than changes
between versions. Second, dynamic detectors require the package to be executed,
whereas this study restricts itself to static analysis.

\subsection{Malicious update detection}

This line of work takes the change between versions as the unit of analysis.
RogueOne flags new data flows that cross trust domains between the old and the
new version. The approach detects 75.6\% of rogue updates with 4.7\% false
positives on a dataset of single updates \citep{sofaer2024rogueone}. RogueOne's
analysis also reveals the blind spot of differential approaches. When the
malicious payload only uses APIs that the previous version already used, no new
data flow appears, and its default configuration therefore misses all 10 cases of
disguised attacks in that category \citep{sofaer2024rogueone}.

\citet{froh2023differential} express suspicious behavior as CodeQL queries and
compare query results across versions. The approach identifies all malicious
versions in a dataset of packages with a history of malicious code, with an
average of 1.3\% false alarms on popular legitimate npm packages.

\citet{wyss2025evaluating} bring the same unit of analysis to LLMs through the
smart-diff. Their analysis shows that false positives stem mainly from sensitive
operations, large-scale code removal (including packages replaced by a deprecation
notice), and dependency changes. A combination of light data flow obfuscation and
payload splitting can also cause false negatives.

\subsection{LLMs for malicious package detection}

SocketAI shows that GPT-4 outperforms CodeQL-based static analysis, with a 16\%
improvement in precision and a 9\% improvement in F1, on 5,115 npm packages of
which 2,180 are malicious \citep{zahan2025leveraging}. MalPacDetector uses an LLM
to generate features automatically instead of relying on expert-defined features.
The result is a 1.3\% false positive rate and a 7.5\% false negative rate, along
with 39 new malicious packages confirmed by the npm security team
\citep{wang2025malpacdetector}.

IntelGuard retrieves semantically similar malicious examples from a knowledge
base of 4,420 malicious code contexts and then uses GPT-4.1-mini to assess
whether the code's behavior matches the package's function
\citep{guo2026bridging}. In the benchmark of \citet{guo2026effective}, this
approach of grounding the assessment in retrieved evidence achieves an F1 of
95.98\%. By contrast, SocketAI, which asks the model to assess the code without
additional evidence, reaches a recall of only 44.41\%.

For open-weight models, \citet{nguyen2025taint} slice code along taint flows
toward sensitive APIs before analyzing it with DeepSeek-Coder 6.7B. The approach
achieves 87.04\% accuracy on 2,537 packages. In the PyPI ecosystem, an evaluation
of 13 LLMs finds that the models recognize common code patterns but miss
indicators that require broader context or an understanding of the author's
intent \citep{ryan2026evaluation}.

\subsection{Retrieval and examples in prompts for security analysis}

Vulnerability detection research offers two lessons relevant to the design of
contrastive examples. The first lesson: distinguishing faulty code from similar
but safe code is a hard task for LLMs. Vul-RAG shows that LLMs achieve an
accuracy of only 0.06--0.14 in distinguishing vulnerable code from similar
patched code. With multi-dimensional vulnerability knowledge, namely functional
semantics, root cause, and fixing approach, accuracy improves by 16--24\%
\citep{du2026vulrag}. PrimeVul demonstrates that insufficiently realistic
benchmarks overstate performance: a 7B model that achieves 68.26\% F1 on BigVul
achieves only 3.09\% on PrimeVul \citep{ding2025vulnerability}.

The second lesson: the benefit of retrieval for open-weight models cannot be
assumed. A replication of Vul-RAG with a range of open-weight models finds that
pairwise accuracy stalls around 0.30, and that increasing model capacity alone
does not improve it much \citep{kaniewski2026revisiting}.

Outside the security domain, the in-context learning literature shows that the
selection, ordering, and labels of examples affect predictions
\citep{liu2022makes,zhao2021calibrate,min2022rethinking}. In addition, positive
and negative examples have been proposed as a way to clarify the desired output
\citep{gao2024customizing}. These findings motivate our decision to control the
number and ordering of examples rather than assuming that retrieval necessarily
helps.

\subsection{Description--behavior consistency}

The idea of checking whether a program's behavior matches its description
developed in the mobile application ecosystem. WHYPER uses natural language
processing to find sentences in an application's description that explain why the
application needs a given permission \citep{pandita2013whyper}. AutoCog measures
the consistency between descriptions and permissions. Across eleven permissions,
the system achieves an average precision of 92.6\% and a recall of 92.0\%, and
was then applied to 45,811 applications \citep{qu2014autocog}. CHABADA clusters
more than 22,500 Android applications by the topics of their descriptions and
then flags outlier API usage within each cluster \citep{gorla2014checking}. For
npm packages, IntelGuard applies similar reasoning implicitly
\citep{guo2026bridging}. No study, however, has measured the effect of an
explicit intent profile separately from the other components.

\subsection{Position of this study}

Table~\ref{tab:positioning} summarizes the position of IntentDiff-RAG relative to
the closest prior work. This study adopts the update-level unit of analysis from
RogueOne, \citet{froh2023differential}, and \citet{wyss2025evaluating}, and adopts
the ideas of retrieval and function-based reasoning from IntelGuard. It differs
in three aspects that are tested under controlled conditions:
\begin{enumerate}
  \item the use of open-weight models of 3--4 billion parameters at the update
  level;
  \item contrastive legitimate examples paired with malicious examples; and
  \item an explicit intent profile drawn from the previous version.
\end{enumerate}
Following the lessons from Vul-RAG, PrimeVul, and its replication, the evaluation
is carried out on hard negatives, that is, legitimate updates that add sensitive
capabilities, and not only on typical updates.

\begin{table*}[t]
\caption{Position of IntentDiff-RAG relative to the closest prior work. A ``--''
means that the component is not used or not described in the sources we examined.
The \emph{Retrieval} column refers to retrieving examples into the LLM prompt.}
\label{tab:positioning}
\small
\begin{tabular}{p{3.2cm} p{1.6cm} p{3.9cm} p{3.2cm} p{2.4cm} p{1.8cm}}
\hline
Study & Unit of analysis & Approach & Language model & Retrieval & Intent \\
\hline
MalOSS \citep{duan2021towards} & Package & Metadata, static, dynamic & -- & -- & -- \\
Amalfi \citep{sejfia2022practical} & Package & Classifier, reproduction, clone detection & -- & -- & -- \\
DONAPI \citep{huang2024donapi} & Package & Static and dynamic API sequences & -- & -- & -- \\
OSCAR \citep{zheng2024towards} & Package & Dynamic (sandbox, fuzzing) & -- & -- & -- \\
\citet{froh2023differential} & Update & Differential static analysis (CodeQL) & -- & -- & -- \\
RogueOne \citep{sofaer2024rogueone} & Update & Differential data flow & -- & -- & -- \\
SocketAI \citep{zahan2025leveraging} & Package & LLM & Closed (GPT-3, GPT-4) & -- & -- \\
\citet{wyss2025evaluating} & Update & Smart-diff + LLM & Closed (GPT-4o, GPT-4-Turbo, Gemini 2.0 Flash) & -- & -- \\
IntelGuard \citep{guo2026bridging} & Package & RAG + LLM & Closed (GPT-4.1-mini) & Malicious examples & Implicit \\
\citet{nguyen2025taint} & Package & Taint slicing + LLM & Open-weight (6.7B) & -- & -- \\
\textbf{IntentDiff-RAG} & \textbf{Update} & \textbf{Smart-diff + RAG + LLM} & \textbf{Open-weight (3--4B)} & \textbf{Malicious + contrastive legitimate} & \textbf{Explicit} \\
\hline
\end{tabular}
\end{table*}

\section{Method}\label{sec:method}

\section{Experimental setup}\label{sec:setup}

\section{Results}\label{sec:results}

\section{Discussion}\label{sec:discussion}

\section{Conclusion}\label{sec:conclusion}

\printcredits

\bibliographystyle{cas-model2-names}
\bibliography{cas-refs}

\end{document}